%%=============================================================================
%% Voorwoord
%%=============================================================================

\chapter*{\IfLanguageName{dutch}{Woord vooraf}{Preface}}%
\label{ch:voorwoord}

%% TODO:
%% Het voorwoord is het enige deel van de graduaatsproef waar je vanuit je
%% eigen standpunt (``ik-vorm'') mag schrijven. Je kan hier bv. motiveren
%% waarom jij het onderwerp wil bespreken.
%% Vergeet ook niet te bedanken wie je geholpen/gesteund/... heeft

Voor u ligt de graduaatsproef "Automatisatie van het Release Verificatieproces". Dit rapport vormt het sluitstuk van mijn opleiding Graduaat Programmeren aan Hogent.

Van februari tot en met juni 2026 heb ik met veel enthousiasme stage gelopen bij het softwarebedrijf Amaron. Tijdens deze periode kreeg ik de kans om me te verdiepen in de complexe, maar fascinerende wereld van DevOps en release management. De opdracht om een geautomatiseerd verificatieplatform te ontwikkelen voor hun groeiende ecosysteem van componenten was een behowelijke technische uitdaging, maar bovenal een enorm leerzaam proces. Ik heb niet alleen mijn technische vaardigheden rondom API's en automatisering sterk kunnen verbeteren, maar ook geleerd hoe cruciaal goede communicatie en gestroomlijnde workflows zijn binnen een groot ontwikkelteam.

Dit eindwerk was echter niet tot stand gekomen zonder de hulp, begeleiding en steun van een aantal mensen die ik hier graag in de bloemetjes wil zetten.

In de eerste plaats wil ik mijn stagebegeleider bij Amaron, Koen Casier, bedanken. Bedankt voor de tijd, de technische inzichten en het vertrouwen dat ik kreeg om zelfstandig aan dit project te werken. Daarnaast wil ik ook het volledige DevOps-team en de developers van Amaron bedanken voor de warme ontvangst, de fijne samenwerking en het beantwoorden van mijn vele vragen.

Ik hoop dat u dit rapport met evenveel plezier leest als ik heb gehad bij het realiseren van dit project.